\section{Completion of the three-prime classification}
\label{sec:three-prime-classification}

The largest quotient root supplies the final obstruction. Its open unit
interval completes endpoint one without deciding the quartic's smallest
root or enumerating integer exponent points on the endpoint surface.

\begin{theorem}
\label{thm:largest-root-unit-interval}
Let $3\leq a\leq b\leq c$ be real, with $b-a\geq3$ and $c\geq a^2-a$.
Put $s=a+b+c$ and $M=bc+s$. Then
\[
 M<\lambda_{\max}(C)<M+1.
\]
For integer exponents in this domain, the ideal intersection graph is
Laplacian nonintegral. The real interval has a written proof with no
finite exponent base.
\end{theorem}

\begin{proof}
The quotient determinant gives
\[
 h_C(M)=-a^2bc\,T(a,b,c),
\]
where
\[
 T=(2b^2-ab+b-a)c^2-[(a-1)b^2+a^2]c-ab(a+b).
\]
Put $P(a,b,c)=h_C(M+1)$. Under
$a=3+m$, $b=a+3+u$, $c=a^2-a+v$, with $m,u,v\geq0$,
the complete identities in Appendix~\ref{sec:largest-root-identities}
express $T$ and $P$ as polynomials with respectively $36$ and $170$
strictly positive coefficients, including constants $1404$ and $513220$.
Thus $h_C(M)<0<h_C(M+1)$, giving a root in $(M,M+1)$.

To identify it as the largest root, use the symmetric similarity $S$
from Theorem~\ref{thm:endpoint-one-spectrum}. For $x>c$, the pair-support
block of $xI-S$ is positive definite. Its Schur complement is
\[
 H(x)=\operatorname{diag}(\ell_a,\ell_b,\ell_c)+zz^T,\qquad
 z=(\sqrt a,\sqrt b,\sqrt c)^T,\qquad
 \ell_t=x-s-\frac{xabc}{t(x-t)}.
\]
At $x=M+1$, both $x-b$ and $x-c$ exceed $bc$. Hence
\begin{align*}
 \ell_b&=c(b-a)+1-\frac{abc}{x-b}>c(b-a)+1-a>0,\\
 \ell_c&=b(c-a)+1-\frac{abc}{x-c}>b(c-a)+1-a>0.
\end{align*}
Here $c\geq a^2-a\geq a+3$. The $b,c$ principal block of $H$ is
positive definite, and $\det(xI-S)=xP>0$ gives $\det H>0$.
Its remaining scalar Schur complement is therefore positive.
Consequently $xI-S$ is positive definite, and every eigenvalue is less
than $M+1$. Since $\det(MI-S)=Mh_C(M)<0$, not all eigenvalues can
lie below $M$, so the largest lies in the stated interval. For integer exponents,
the established complement and universal-class lifts transfer its
nonintegrality to the ideal intersection graph.
\end{proof}

The first two middle gaps use exponent-root brackets rather than the
largest-root interval. The following proof records their signs explicitly
so that the classification does not rely on an unstated uniqueness claim.

\begin{lemma}
\label{lem:endpoint-one-small-gaps}
Let $a\geq2$ be real, $d\in\{1,2\}$ and $b=a+d$. Put
$L_1=2a^2-a-2$ and $L_2=2a^2+a-6$. As a polynomial in $c$,
$h_C(1)$ has exactly one root above $b$, lying in $(L_d,L_d+1)$.
Consequently, for integer $a\geq2$, no integer $c>b$ satisfies $h_C(1)=0$.
\end{lemma}
\begin{proof}
For fixed $a,b$, write $Q(c)=h_C(1)=A c^3+B c^2+D c+E$, where
\begin{align*}
 A&=a^2+b^2-1,\\
 B&=a^3-4a^2b^2+2a^2b-a^2+2ab^2+ab-3a+b^3-b^2-3b+3,\\
 D&=(a-1)(b-1)(2ab+3a+3b-3),\\
 E&=(a-1)(b-1)(a+b-1)(ab+a+b-1).
\end{align*}
Both $A$ and $E$ are positive. Put $L_1=2a^2-a-2$ and
$L_2=2a^2+a-6$. Direct expansion gives
\begin{align*}
 -Q(L_1)&=4a(a+1)(a^4-4a^2-2a+6),\\
 Q(L_1+1)&=4a(a-1)(a+1)(a^3-a^2+1),\\
 -Q(L_2)&=16a^5+72a^4-76a^3-433a^2+69a+595,\\
 Q(L_2+1)&=4(a+2)(2a^5-2a^4-17a^3+25a^2+28a-42),
\end{align*}
where each pair uses its own $b=a+d$. The complete positive vectors
at $a=2+m$, including those for $-Q(b)$, are in
Appendix~\ref{sec:largest-root-identities}. Thus $Q(0)>0>Q(b)$ and
$Q(L_d)<0<Q(L_d+1)$. Moreover,
\[
 L_1-b=1+6(a-2)+2(a-2)^2>0,\qquad
 L_2-b=2(a-2)(a+2)\geq0.
\]
The second equality occurs only at $a=2$; the final open interval
still lies strictly above $b$ there.
The cubic has one root in each of $(-\infty,0)$, $(0,b)$ and
$(L_d,L_d+1)$. These three distinct roots exhaust its degree.
The unique root above $b$ is therefore in $(L_d,L_d+1)$.
For integer $a$, the endpoints are consecutive integers, proving the
nonvanishing conclusion without a finite base.
\end{proof}

\begin{theorem}[Three-prime nonintegrality]
\label{thm:three-prime-classification}
For distinct primes $p,q,r$ and positive integers $a,b,c$, the ideal
intersection graph of $\mathbb Z_{p^a q^b r^c}$ is Laplacian nonintegral.
This global conclusion combines written proofs with the established
small-minimum and endpoint-two finite certificates.
\end{theorem}
\begin{proof}
Sort the exponents. Repeated entries are covered by
Theorem~\ref{thm:repeated}; minima at most seven by
Corollary~\ref{cor:minimum-seven}. Suppose $8\leq a<b<c$ and the
spectrum is integral. Theorem~\ref{thm:uniform-low-root} supplies a
positive quotient root in $(0,3)$, necessarily one or two.
Theorem~\ref{thm:endpoint-two-completion} excludes the second alternative,
so $h_C(1)=0$.

For $b-a=1$ or $2$, Lemma~\ref{lem:endpoint-one-small-gaps} contradicts
integer $c>b$. Otherwise $b-a\geq3$, and
Theorem~\ref{thm:endpoint-one-spectrum} gives
$b+c\geq2a^2-2a+1$. Since $c\geq b$, we have $c>a^2-a$.
Theorem~\ref{thm:largest-root-unit-interval} supplies a noninteger
quotient root, again contradicting integrality. These cases exhaust
all positive exponent triples.

The minimum-through-seven theorem retains its $27\,562$-triple base
and earlier small-minimum dependencies; the endpoint-two theorem
retains its $8\,658$-triple base. The new largest-root and small-gap
steps have written unbounded proofs. No additional exponent scan or
integer-point enumeration is used in this completion.
\end{proof}
